\newtoggle{withbonus}
\settoggle{withbonus}{false} % set true to show bonus points in grade table

% setting underlying course name
\newcommand{\mycourse}{Informatik}

% name of the class
\newcommand{\myclass}{\faPeace\faIceCream~Klasse 2f~\faIceCream\faPeace}

% set date
\newcommand{\mydate}{Donnerstag, 18. Juni 2026}

\title{\textbf{randomisierte Algorithmen, Graphen, Datenintegrität}}
\author{Thomas Graf\\
\mycourse, \myclass}
\date{\mydate}
\maketitle
\thispagestyle{headandfoot}

\begin{center}
	\fbox{\parbox{140mm}{\centering
			\begin{itemize}
				\item Dauer der Prüfung: $60$ Minuten
				\item erlaubte Hilfsmittel:
				\begin{itemize}
					\item Schreibzeug
					\item \textbf{ohne Taschenrechner}
				\end{itemize}
			\end{itemize}}}
\end{center}

\vspace{10mm}

\begin{center}
	Vorname: \rule{45mm}{0.8pt}\qquad Nachname: \rule{45mm}{0.8pt}
\end{center}

\vspace{20mm}

\begin{center}
	\iftoggle{withbonus}{
		\combinedgradetable[h][questions]
	}{
		\gradetable[h][questions]
	}
\end{center}

\vspace{15mm}

\begin{center}
	Note:
	\vspace{10mm}

	\rule{8mm}{1.2pt}
\end{center}

\clearpage

\begin{questions}

% 2c
\question
\begin{parts}
	\part[1] Bestimmen Sie $\text{Nummer}(0101111)$. \vspace{20mm}

	\part[1] Bestimmen Sie ein binäres Wort (einen binären String) $x$, sodass $\text{Nummer}(x) = 57$. \vspace{20mm}

	\part[1] Berechnen Sie $\abs{\text{Bin}(3099)}$. \vspace{20mm}

	\part
	Sei $y$ die binäre Darstellung ($2$-adische Darstellung) der drittgrössten Zahl, welche mit $7$ Bits dargestellt werden kann.
	\begin{subparts}
		\subpart[1] Wie sieht $y$ aus? \vspace{15mm}
		\subpart[1] Bestimmen Sie $\text{Nummer}\lr{y}$. \vspace{15mm}
	\end{subparts}

	\part[1] Berechnen Sie die Zahl $\ceil{\log_6\lr{180}}$. \vspace{20mm}

	\part[1] Bestimmen Sie $\text{Prim}(37)$. \vspace{20mm}

	\part[1] Nennen Sie positive natürliche Zahlen $x, y$ und $z$, sodass zwar
	\begin{align*}
		x < y < z
	\end{align*}
	gilt, aber nicht
	\begin{align*}
		\abs{\text{Bin}(x)} < \abs{\text{Bin}(y)} < \abs{\text{Bin}(z)}.
	\end{align*}
	\vspace{20mm}
\end{parts}
\droptotalpoints

\begin{solution}
\begin{enumerate}
	\item \(\text{Nummer}(0101111) = 32 + 8 + 4 + 2 + 1 = 47\).

	\item \(57 = 32 + 16 + 8 + 1 = 2^5 + 2^4 + 2^3 + 2^0\). Also ist \(x = 111001\).

	\item Die Zahl \(3099\) liegt zwischen \(2^{11} = 2048\) und \(2^{12} = 4096\). Daher gilt
	\[
		\abs{\text{Bin}(3099)} = 12.
	\]

	\item
	\begin{enumerate}
		\item Die drittgrösste Zahl mit \(7\) Bits ist \(125 = 1111101_2\). Also ist \(y = 1111101\).
		\item \(\text{Nummer}(y) = 125\).
	\end{enumerate}

	\item Es gilt \(6^2 = 36 < 180 < 216 = 6^3\). Daher ist
	\[
		\ceil{\log_6\lr{180}} = 3.
	\]

	\item Die \(37\). Primzahl ist \(157\).

	\item Ein mögliches Beispiel ist \(x = 3\), \(y = 4\), \(z = 5\). Dann gilt
	\[
		\abs{\text{Bin}(3)} = 2, \qquad
		\abs{\text{Bin}(4)} = 3, \qquad
		\abs{\text{Bin}(5)} = 3.
	\]
	Also gilt zwar \(3 < 4 < 5\), aber nicht \(2 < 3 < 3\).
\end{enumerate}
\end{solution}

\question[3]
Gegeben sind die beiden Datensätze
\begin{itemize}
	\item $x := 100110$ in Zürich und
	\item $y := 010111$ in Boston.
\end{itemize}
Berechnen Sie die exakte Fehlerwahrscheinlichkeit des Kommunikationsprotokolls in diesem Fall.
\vspace{40mm}

\begin{solution}
	Die Datensätze haben Länge \(n = 6\). Also wird eine Primzahl \(p \leq n^2 = 36\) zufällig gewählt.

	Es gilt
	\[
		100110_2 = 38
		\qquad \text{und} \qquad
		010111_2 = 23.
	\]
	Damit ist
	\[
		\abs{x-y} = 15.
	\]
	Ein Fehler tritt genau dann auf, wenn \(p\) die Zahl \(15\) teilt.

	Die Primzahlen kleiner oder gleich \(36\) sind
	\[
		2,3,5,7,11,13,17,19,23,29,31.
	\]
	Davon teilen nur \(3\) und \(5\) die Zahl \(15\). Somit beträgt die exakte Fehlerwahrscheinlichkeit
	\[
		\frac{2}{11}.
	\]
\end{solution}

\clearpage

\question
\begin{parts}
	\part[2]
	Für die Fehlerwahrscheinlichkeit $r$ des Protokolls haben wir die Abschätzung
	\begin{align*}
		r \leq \frac{n-1}{\text{Prim}(n^2)}
		\leq \frac{n-1}{n^2/\ln(n^2)}
		\leq \frac{\ln(n^2)}{n}
	\end{align*}
	gefunden. Dabei ist $n$ die Länge der jeweiligen Datensätze in Bits. Erklären Sie, warum wir im Protokoll eine Primzahl $p \leq n^2$ wählen und nicht lediglich eine Primzahl $p \leq n$.
	\vspace{40mm}

	\part[1]
	Neben der wiederholten Durchführung des Protokolls könnte man die Fehlerwahrscheinlichkeit auch verringern, indem man anstelle unseres Protokolls $R$ ein leicht angepasstes Protokoll $R'$ betrachtet. $R'$ arbeitet exakt so wie $R$ mit der einzigen Änderung, dass $R'$ eine Primzahl zufällig aus der Menge $\lrc{2, 3, \ldots , n^r}$ statt aus der Menge $\lrc{2, 3, \ldots , n^2}$ wählt. Berechnen Sie die Fehlerwahrscheinlichkeit von $R'$ für jede natürliche Zahl $r > 2$ und vergleichen Sie diese mit der Fehlerwahrscheinlichkeit von $R$.
	\vspace{40mm}
\end{parts}
\droptotalpoints

\question
\textbf{Graphentheorie: Eulerwege}

Gegeben ist ein ungerichteter Graph mit den Knoten
\[
	V = \{A, B, C, D, E\}
\]
und den Kanten
\[
	E = \{\{A,B\}, \{A,C\}, \{B,C\}, \{B,D\}, \{B,E\}, \{C,D\}, \{C,E\}, \{D,E\}\}.
\]

\begin{parts}
	\part[1] Bestimmen Sie die Grade der Knoten $A$ und $B$. \vspace{20mm}
	\part[1] Existiert in diesem Graphen ein Eulerkreis? Begründen Sie Ihre Antwort kurz mit Hilfe der Knotengrade. \vspace{20mm}
\end{parts}
\droptotalpoints

\begin{solution}
\begin{enumerate}
	\item Grad($A$) = 2 und Grad($B$) = 4.

	\item Ein Eulerkreis existiert genau dann, wenn alle Knotengrade gerade sind. Es gilt
	\[
		\deg(A)=2, \quad \deg(B)=4, \quad \deg(C)=4, \quad \deg(D)=3, \quad \deg(E)=3.
	\]
	Da \(D\) und \(E\) ungeraden Grad haben, existiert kein Eulerkreis.
\end{enumerate}
\end{solution}

\clearpage

\question[3]
\textbf{Modifikation des Dijkstra-Algorithmus}

\begin{lstlisting}[language=Python,numbers=none,escapeinside={(*}{*)},caption=\texttt{Dijkstra-Algorithmus},label=listing:Dijkstra]
### Dijkstra-Algorithmus ###

## Initialisierung ##
für alle Knoten u des Graphen setze:
    (*$d_s[u] = \infty$*)
(*$d_s[s] \leftarrow 0$*)
(*$S \leftarrow \emptyset$*)
(*$Q \leftarrow $*) Menge aller Knoten des Graphen

## Schritte ##
solange (*$Q$*) nicht leer ist, mach:
    (*$u\leftarrow$*) Knoten aus (*$Q$*) mit kleinstem (*$d_s[u]$*)
    (*$Q \leftarrow Q \setminus \lrc{u}$*)
    (*$S\leftarrow S \cup \lrc{u}$*)

    für alle von (*$u$*) direkt erreichbaren Knoten (*$v$*) mach:
        # ist der Weg von (*$s$*) nach (*$v$*) kürzer über (*$u$*)?
        falls (*$d_s[u] + c(u,v) < d_s[v]$*) dann:
            (*$d_s[v] \leftarrow d_s[u] + c(u,v)$*)
\end{lstlisting}

Im Standard-Dijkstra-Algorithmus siehe \cref{listing:Dijkstra} suchen wir den Pfad mit der minimalen \textit{Summe} der Kantengewichte. Nehmen wir nun an, wir suchen in einem Kommunikationsnetzwerk den Weg mit der \textbf{maximalen Zuverlässigkeit}.

Jeder Kante $(u, v)$ ist eine Wahrscheinlichkeit $P((u, v)) \in \ioc{0}{1}$ zugeordnet, dass die Verbindung funktioniert. Die Zuverlässigkeit eines Pfades ist das \textbf{Produkt} der Wahrscheinlichkeiten aller Kanten auf diesem Pfad.

Kreuzen Sie für die drei wesentlichen Änderungen am Pseudocode jeweils die korrekte Option an, damit der Algorithmus den zuverlässigsten Pfad findet:

\textbf{1. Initialisierung:}
\begin{checkboxes}
	\choice $d_s[u] = \infty$ und $d_s[s] = 0$
	\CorrectChoice $d_s[u] = 0$ und $d_s[s] = 1$
	\choice $d_s[u] = 1$ und $d_s[s] = 0$
\end{checkboxes}

\textbf{2. Knotenauswahl:}
\begin{checkboxes}
	\choice Wähle $u$ mit dem kleinsten $d_s[u]$
	\choice Wähle $u$ zufällig aus $Q$
	\CorrectChoice Wähle $u$ mit dem grössten $d_s[u]$
\end{checkboxes}

\textbf{3. Relaxation:}
\begin{checkboxes}
	\choice Falls $d_s[u] + P((u, v)) < d_s[v]$
	\CorrectChoice Falls $d_s[u] \cdot P((u, v)) > d_s[v]$
	\choice Falls $d_s[u] \cdot P((u, v)) < d_s[v]$
\end{checkboxes}

\begin{solution}
\begin{enumerate}
	\item Alle Distanzwerte \(d_s[u]\) werden zu Beginn auf \(0\) gesetzt. Der Startknoten erhält den Wert \(d_s[s] = 1\).

	\item In jedem Schritt muss der Knoten \(u \in Q\) mit dem grössten aktuellen Wert \(d_s[u]\) gewählt werden.

	\item Der Vergleich und das Update verwenden Multiplikation:
	\[
		d_s[u] \cdot P((u,v)) > d_s[v].
	\]
\end{enumerate}
\end{solution}

\clearpage

\question[2]
\textbf{EAN-13 Prüfziffer}

Ein EAN-13-Code besteht aus 12 Ziffern $x_1 \ldots x_{12}$ und einer Prüfziffer $x_{13}$. Die gewichtete Summe $s$ berechnet sich nach der Formel:
\[
	s = x_1 + x_3 + x_5 + x_7 + x_9 + x_{11}
	+ 3 \cdot (x_2 + x_4 + x_6 + x_8 + x_{10} + x_{12}).
\]
Die Prüfziffer $x_{13}$ wird anschliessend so gewählt, dass $s + x_{13}$ durch 10 teilbar ist.

Berechnen Sie die Prüfziffer $x_{13}$ für die folgende Ziffernfolge:
\[
	\texttt{541230981246}.
\]
\vspace{30mm}

\begin{solution}
	Summe der ungeraden Stellen:
	\[
		5 + 1 + 3 + 9 + 1 + 4 = 23.
	\]
	Summe der geraden Stellen:
	\[
		4 + 2 + 0 + 8 + 2 + 6 = 22.
	\]
	Diese wird mit \(3\) gewichtet:
	\[
		3 \cdot 22 = 66.
	\]
	Die Gesamtsumme ist
	\[
		s = 23 + 66 = 89.
	\]
	Damit \(89 + x_{13}\) durch \(10\) teilbar ist, muss \(x_{13} = 1\) sein.
\end{solution}

\question
\textbf{Fehlererkennung und -korrektur}

Gegeben ist eine binäre Kodierung mit den folgenden drei gültigen Codewörtern:
\begin{itemize}
	\item $C_1 = 101010$
	\item $C_2 = 110110$
	\item $C_3 = 001011$
\end{itemize}

\begin{parts}
	\part[1] Bestimmen Sie den minimalen Abstand dieser Kodierung. \vspace{15mm}
	\part[1] Wie viele Bitfehler können mit dieser Kodierung \textit{sicher erkannt} werden? \vspace{15mm}
	\part[1] Wie viele Bitfehler können mit dieser Kodierung \textit{sicher korrigiert} werden? \vspace{15mm}
\end{parts}

\begin{solution}
\begin{enumerate}
	\item Die Abstände sind
	\[
		d(C_1, C_2) = 3, \qquad
		d(C_1, C_3) = 2, \qquad
		d(C_2, C_3) = 3.
	\]
	Der minimale Abstand ist also \(d_{\min} = 2\).

	\item Sicher erkannt werden können
	\[
		d_{\min} - 1 = 1
	\]
	Bitfehler.

	\item Sicher korrigiert werden können
	\[
		\left\lfloor \frac{d_{\min} - 1}{2} \right\rfloor = 0
	\]
	Bitfehler.
\end{enumerate}
\end{solution}

\end{questions}